% =====================================================================
%  DRAFT — AI-ASSISTED SCAFFOLD, FOR GROUP REVIEW ONLY.
%  The executable assurance case (extends 3.6) + Chapter 4 Critical Evaluation.
%  NOT for verbatim submission: re-voice in your own words and verify every
%  figure against the code/assurance-case.md before the report is submitted.
%  Facts from branch clarence/c-block-05 (validate_assurance.py: graph SOUND,
%  15 claims, 0 authorised) 2026-08-19. Requires booktabs; Table 5 = tab:assurance.
% =====================================================================

\subsection{The Executable Assurance Case}
\label{sec:assurance}
The controls described in the preceding sections are not asserted individually but are consolidated
into an executable assurance case, which is the mechanism by which the project reasons about its own
maturity. Each claim is represented as a small graph linking the promise to the affected party and the
harm it addresses, to the control that keeps it, to the executable test or tests that evidence that
control, and finally to the two human gates it still requires: a named non-author reviewer and a named
authority. A validator runs every linked test live and additionally performs orphan detection, so that
a claim missing any element---an unowned control, an unbacked test, or an unnamed reviewer or
authority---cannot silently pass. The case presently holds fifteen such claims (Table~\ref{tab:assurance}).
Its dependency graph is sound and every linked test is green, yet the maturity verdict of all fifteen is
\textsc{blocked}, because the two human gates are, by design, still open. This is a deliberate design
decision rather than an incidental shortcoming. The project's central thesis---that a system must
distinguish what its evidence supports from what it cannot determine, and must never convert the absence
of evidence into a positive claim---is here turned upon the project's own governance: an assurance
dashboard that reported ``not yet reviewed'' as ``authorised'' would commit precisely the error the
applications are built to prevent. The honest, and defensible, state is therefore that the executable
evidence is complete while the human authorisation remains outstanding.

\section{Critical Evaluation}
\label{ch:eval}
The contribution of this stream should be read as a research demonstration of a proposed design, and
its limits are best stated in the same disciplined terms as its claims. What has been demonstrated is
that the applications and their supporting contracts satisfy a set of executable invariants: the public
and staff surfaces remain evidence-symmetric and value-level non-interfering, unknown and conflicting
states are never rendered as negative facts, unsupported model text cannot reach a user, and governed
staff actions cannot skip review. What has not been demonstrated is any human-facing effect. None of the
fifteen assurance claims has been authorised; each remains blocked pending a named non-author review and
an authority decision that the author is explicitly not entitled to issue. The distinction is not
pedantic: ``the evidence is green'' is a statement about the code, whereas ``the claim is authorised''
is a statement about a person with standing having accepted the residual risk, and only the second
licenses a maturity claim.

The design was specified with explicit falsification conditions, and it is instructive to evaluate the
work against them. The contribution would narrow or fail if the public and staff surfaces disagreed on
the retained semantics, if restricted staff information leaked into the public plane, if the interface
bypassed model abstention or the certified order, or if unsupported model content reached users; the
executable gates target each of these, and none is currently violated on the tested fixtures. That the
gates are not merely decorative was borne out during review, when a disclosure defect---the public
service-failure page named an internal pipeline stage held in a field that had been misclassified as
public-safe---was detected, reclassified to a staff-only class, and closed with a regression test.
Reporting this openly is deliberate: it shows the control functioning as intended rather than concealing
a lapse, and a single such finding strengthens rather than weakens the assurance argument, because it
demonstrates that the case is capable of failing.

Several limits are load-bearing and should not be understated. The accessibility position is
deliberately scoped as ``evaluated toward WCAG~2.2~AA within a tested matrix'' rather than as
conformance, because the automated checks performed are only one input and the manual keyboard,
screen-reader and assistive-technology testing, together with an independent accessibility review,
remain outstanding; a single examiner-found failure would falsify a blanket conformance claim, whereas
the scoped statement survives \cite{power2012guidelines, w3c2023wcag22}. The evaluation manifest,
similarly, currently carries three tasks that are adequate as demonstration fixtures for the
confound-proof but are not a benchmark task set: an independent review of the manifest against the
evaluation estimands established that, at three tasks, the crossed-design power ceiling sits at
approximately $0.18$, so that no realistic participant count reaches a usable target, and a benchmark
would require on the order of twenty or more tasks. The staff role gate is a demonstration stub rather
than production identity management, the conversational parser is deterministic and stands in for a
language model that would sit behind the same contract, and the data-protection and research-ethics
determinations that would license any human-facing study are institutional decisions that are
progressed and evidenced but cannot be self-certified. Each of these is recorded rather than obscured,
and each has a defined route to closure.

Taken together, these limits define the distance between the present research demonstration and a
system that could make stronger claims. Closing that distance is a matter of completing the human half
of the assurance case---the non-author reviews, the named authorities, and the team ratifications---and
of executing the evaluation the instruments were built to support: minting the benchmark task set
against a stated power target, and running the manual accessibility and adversarial security batteries.
The contribution of the work is therefore not a claim of effectiveness, which the evaluation is not yet
positioned to make, but the demonstrated integration of an evidence-symmetric, authority-asymmetric
application pair over a single certified contract, together with the executable discipline that keeps
its claims honest---a discipline whose value is precisely that it renders the outstanding work visible
and auditable rather than assumed.
